\documentclass[a4paper,11pt]{article}
\usepackage{lmodern}
\usepackage[T1]{fontenc}
\usepackage[utf8]{inputenc}
\usepackage{amsmath,amssymb,amsthm}
\usepackage{longtable,booktabs,array,calc}
\usepackage{hyperref}
\usepackage[margin=25mm]{geometry}

\newtheorem{theorem}{Theorem}[section]
\newtheorem{proposition}[theorem]{Proposition}
\newtheorem{lemma}[theorem]{Lemma}
\newtheorem{corollary}[theorem]{Corollary}
\newtheorem{definition}[theorem]{Definition}
\newtheorem{remark}[theorem]{Remark}
\newtheorem{observation}[theorem]{Observation}

\providecommand{\tightlist}{\setlength{\itemsep}{0pt}\setlength{\parskip}{0pt}}
\newcounter{none}

\begin{document}

\section{Geometric Classification of Spin Derived from the Orientation
Structure of the 5-Dimensional
Orthoplex}\label{geometric-classification-of-spin-derived-from-the-orientation-structure-of-the-5-dimensional-orthoplex}

\subsection{--- A Unified Derivation of Spin Values, Statistics, and
Force Directionality from Integer Geometry
---}\label{a-unified-derivation-of-spin-values-statistics-and-force-directionality-from-integer-geometry}

\textbf{Author:} Noriaki Kihara (WF System Co., Ltd.)

\textbf{Date:} April 2026

\textbf{DOI:}
\href{https://doi.org/10.5281/zenodo.19630972}{10.5281/zenodo.19630972}

\textbf{Prerequisite Papers:} -
\href{https://doi.org/10.5281/zenodo.19624957}{Full Spherical Coverage
of Central Projection in Discrete Space and the Number-Theoretic
Necessity of 5-Dimensional Background Space} -
\href{https://doi.org/10.5281/zenodo.19601592}{Formulation of Structural
Requirements for a Theory of Everything} -
\href{https://doi.org/10.5281/zenodo.19604965}{Classification and
Structural Analysis of the 19 Arbitrary Parameters of the Standard
Model}

\begin{center}\rule{0.5\linewidth}{0.5pt}\end{center}

\subsection{0. Purpose of This Paper}\label{purpose-of-this-paper}

In the preceding paper {[}W5{]}, we demonstrated that the 5-dimensional
orthoplex (the 5-dimensional regular cross-polytope) constitutes the
unit cell for dense packing of discrete space, and that under central
projection it is mapped to a hyperrectangle in the 4-dimensional
subjective space. We also established in Proposition 7.1 that the sign
of the volume (orientation) is naturally defined through the determinant
sign of the 5-dimensional edge vectors.

This paper further analyzes the orientation structure and derives the
following:

\begin{enumerate}
\def\labelenumi{\arabic{enumi}.}
\tightlist
\item
  The set of representable spin quantum numbers is limited to six
  values: \(s = 0, \frac{1}{2}, 1, \frac{3}{2}, 2, \frac{5}{2}\), as
  determined by the orientation degrees of freedom of the 5-dimensional
  orthoplex.
\item
  The distinction between half-integer spin (fermions) and integer spin
  (bosons) is derived as a geometric difference between vertex states
  and edge-midpoint states.
\item
  The directionality of forces (attractive only versus attractive and
  repulsive) is determined by the parity of the sign-reversal structure
  of the volume.
\item
  A generational structure of particles naturally emerges from the
  permutation symmetry \(S_3\) of the three spatial axes.
\item
  The sign asymmetry of spin-0 states (10 positive versus 9 negative)
  suggests a geometric origin for matter--antimatter asymmetry.
\end{enumerate}

All results in this paper are stated as propositions of discrete
geometry and number theory, without invoking physical interpretations.
However, correspondences with known particle physics are noted as
remarks.

\begin{center}\rule{0.5\linewidth}{0.5pt}\end{center}

\subsection{1. Fundamental Structure of the 5-Dimensional
Orthoplex}\label{fundamental-structure-of-the-5-dimensional-orthoplex}

\subsubsection{1.1 Definition and Element
Counts}\label{definition-and-element-counts}

The 5-dimensional orthoplex is the regular polytope in \(\mathbb{R}^5\)
with vertices at \(\pm e_i\) (\(i = 1, 2, 3, 4, 5\)) on each coordinate
axis. The edge length is set to 1 (the vertex-to-vertex distance
\(\sqrt{2}\) satisfies \(\ell^2 = 2\), which is an integer).

The number of \(k\)-dimensional faces is given by
\(2^{k+1}\binom{5}{k+1}\).

{\def\LTcaptype{none} % do not increment counter
\begin{longtable}[]{@{}cccc@{}}
\toprule\noalign{}
\(k\) & Element & Count & Computation \\
\midrule\noalign{}
\endhead
\bottomrule\noalign{}
\endlastfoot
0 & Vertices & 10 & \(2^1 \binom{5}{1}\) \\
1 & Edges & 40 & \(2^2 \binom{5}{2}\) \\
2 & Triangular faces & 80 & \(2^3 \binom{5}{3}\) \\
3 & Tetrahedral cells & 80 & \(2^4 \binom{5}{4}\) \\
4 & 5-cell facets & 32 & \(2^5 \binom{5}{5}\) \\
\end{longtable}
}

\subsubsection{1.2 Integrality of Metrics}\label{integrality-of-metrics}

With unit edge length, the metric at each dimension is computed as
follows.

{\def\LTcaptype{none} % do not increment counter
\begin{longtable}[]{@{}
  >{\centering\arraybackslash}p{(\linewidth - 6\tabcolsep) * \real{0.2245}}
  >{\centering\arraybackslash}p{(\linewidth - 6\tabcolsep) * \real{0.1633}}
  >{\centering\arraybackslash}p{(\linewidth - 6\tabcolsep) * \real{0.3878}}
  >{\centering\arraybackslash}p{(\linewidth - 6\tabcolsep) * \real{0.2245}}@{}}
\toprule\noalign{}
\begin{minipage}[b]{\linewidth}\centering
Dimension
\end{minipage} & \begin{minipage}[b]{\linewidth}\centering
Metric
\end{minipage} & \begin{minipage}[b]{\linewidth}\centering
Value per element
\end{minipage} & \begin{minipage}[b]{\linewidth}\centering
Rational?
\end{minipage} \\
\midrule\noalign{}
\endhead
\bottomrule\noalign{}
\endlastfoot
0 & Count & 1 & Rational \\
1 & Edge length & 1 & Rational \\
2 & Area (equilateral triangle) & \(\frac{\sqrt{3}}{4}\) &
\textbf{Irrational} \\
3 & Volume (regular tetrahedron) & \(\frac{\sqrt{2}}{12}\) &
\textbf{Irrational} \\
\end{longtable}
}

Only the 0-dimensional (count) and 1-dimensional (edge length) metrics
are rational. Metrics of dimension 2 and above necessarily involve
irrational numbers. This indicates the privileged status of edges
(1-dimensional elements) in integer geometry.

\subsubsection{1.3 A Numerical Coincidence Specific to 5
Dimensions}\label{a-numerical-coincidence-specific-to-5-dimensions}

The count of 2-dimensional faces (triangles), 80, coincides with the
count of 3-dimensional cells (tetrahedra), 80. This follows from
\(2^3 \binom{5}{3} = 2^4 \binom{5}{4}\) and is specific to 5 dimensions.
It suggests a structure in which face orientation and chirality are in
one-to-one correspondence.

\begin{center}\rule{0.5\linewidth}{0.5pt}\end{center}

\subsection{2. Orientation Degrees of Freedom and Independent Oriented
Subspaces}\label{orientation-degrees-of-freedom-and-independent-oriented-subspaces}

\subsubsection{\texorpdfstring{2.1 Number of Independent Components of
\(k\)-Vectors}{2.1 Number of Independent Components of k-Vectors}}\label{number-of-independent-components-of-k-vectors}

The number of independent oriented \(k\)-dimensional subspaces in
5-dimensional space is \(\binom{5}{k}\).

{\def\LTcaptype{none} % do not increment counter
\begin{longtable}[]{@{}cccc@{}}
\toprule\noalign{}
\(k\) & \(\binom{5}{k}\) & Geometric meaning & Meaning of the sign \\
\midrule\noalign{}
\endhead
\bottomrule\noalign{}
\endlastfoot
0 & 1 & Scalar (existence) & Present / absent \\
1 & 5 & Vector (edge direction) & \(\pm\) direction \\
2 & 10 & Bivector (face orientation) & \(\pm\) rotational sense \\
3 & 10 & Trivector (volume chirality) & \(\pm\) chirality \\
4 & 5 & 4-vector (hypervolume orientation) & \(\pm\) orientation \\
5 & 1 & Pseudoscalar (total orientation) & \(\pm\) total chirality \\
\end{longtable}
}

Total: \(\sum_{k=0}^{5} \binom{5}{k} = 2^5 = 32\)

This value of 32 coincides with the number of 4-dimensional facets
(5-cells) of the 5-dimensional orthoplex.

\subsubsection{2.2 Relation to Pascal's
Triangle}\label{relation-to-pascals-triangle}

\[1, \; 5, \; 10, \; 10, \; 5, \; 1\]

This symmetric sequence reflects the Hodge duality between \(k\) and
\((5-k)\). In particular, the coincidence
\(\binom{5}{2} = \binom{5}{3} = 10\) means that the orientation of
2-dimensional faces and the chirality of 3-dimensional volumes are in
one-to-one correspondence.

\begin{center}\rule{0.5\linewidth}{0.5pt}\end{center}

\subsection{3. Geometric Derivation of Spin
Values}\label{geometric-derivation-of-spin-values}

\subsubsection{3.1 Classification of Axes by Physical
Role}\label{classification-of-axes-by-physical-role}

The five coordinate axes are classified based on symmetry.

\begin{itemize}
\tightlist
\item
  \textbf{Spatial axes}: \(x_1, x_2, x_3\) (three axes, symmetric under
  the permutation group \(S_3\))
\item
  \textbf{Fourth axis}: \(x_4\) (asymmetric with respect to spatial
  axes; first candidate for sign reversal)
\item
  \textbf{Fifth axis}: \(x_5\) (the additional dimension for projection;
  second candidate for sign reversal)
\end{itemize}

The three spatial axes are interchangeable, while the fourth and fifth
axes each play an independent role.

\subsubsection{3.2 Correspondence Between Sign Reversals and
Spin}\label{correspondence-between-sign-reversals-and-spin}

When the coupling matrix \(C\) introduces sign reversals on \(n\)
non-spatial axes, the bosonic spin is \(s = n\).

{\def\LTcaptype{none} % do not increment counter
\begin{longtable}[]{@{}
  >{\centering\arraybackslash}p{(\linewidth - 6\tabcolsep) * \real{0.2941}}
  >{\centering\arraybackslash}p{(\linewidth - 6\tabcolsep) * \real{0.2647}}
  >{\centering\arraybackslash}p{(\linewidth - 6\tabcolsep) * \real{0.2059}}
  >{\centering\arraybackslash}p{(\linewidth - 6\tabcolsep) * \real{0.2353}}@{}}
\toprule\noalign{}
\begin{minipage}[b]{\linewidth}\centering
Sign reversals \(n\)
\end{minipage} & \begin{minipage}[b]{\linewidth}\centering
Structure of \(C\)
\end{minipage} & \begin{minipage}[b]{\linewidth}\centering
Bosonic spin
\end{minipage} & \begin{minipage}[b]{\linewidth}\centering
Fermionic spin
\end{minipage} \\
\midrule\noalign{}
\endhead
\bottomrule\noalign{}
\endlastfoot
0 & \(\mathrm{diag}(+,+,+,+,+)\) & 0 & 1/2 \\
1 & \(\mathrm{diag}(+,+,+,-,+)\) & 1 & 3/2 \\
2 & \(\mathrm{diag}(+,+,+,-,-)\) & 2 & 5/2 \\
\end{longtable}
}

Since there are only two non-spatial axes (\(x_4\) and \(x_5\)), the
maximum number of sign reversals is two.

\textbf{Proposition 3.1.} The spin quantum numbers representable by the
orientation structure of the 5-dimensional orthoplex are limited to the
six values

\[s \in \left\{ 0, \; \frac{1}{2}, \; 1, \; \frac{3}{2}, \; 2, \; \frac{5}{2} \right\}.\]

\emph{Proof.} The three spatial axes do not contribute independent sign
reversals due to their \(S_3\) symmetry. The only axes that can undergo
independent sign reversal are the fourth and fifth axes, totaling two.
Since bosonic spin equals the number of sign reversals, the bosonic
values are \(s = 0, 1, 2\). Since fermionic spin is \(s + \frac{1}{2}\),
the fermionic values are \(s = \frac{1}{2}, \frac{3}{2}, \frac{5}{2}\).
The total is six. \(\square\)

\textbf{Remark.} The experimentally confirmed spins of fundamental
particles are \(0, \frac{1}{2}, 1\) only. The values \(\frac{3}{2}\) and
\(2\) are theoretical predictions (gravitino and graviton,
respectively). All six values fall within the scope of known theoretical
frameworks.

\subsubsection{3.3 Number of Spin States}\label{number-of-spin-states}

A particle of spin \(s\) has \(2s + 1\) values of \(s_z\).

{\def\LTcaptype{none} % do not increment counter
\begin{longtable}[]{@{}cc@{}}
\toprule\noalign{}
\(s\) & Number of states \(2s+1\) \\
\midrule\noalign{}
\endhead
\bottomrule\noalign{}
\endlastfoot
0 & 1 \\
1/2 & 2 \\
1 & 3 \\
3/2 & 4 \\
2 & 5 \\
5/2 & 6 \\
\textbf{Total} & \textbf{21} \\
\end{longtable}
}

\[\sum_{s} (2s+1) = 1 + 2 + 3 + 4 + 5 + 6 = 21 = T(6)\]

21 is the sixth triangular number and also equals \(\binom{7}{2}\).

\begin{center}\rule{0.5\linewidth}{0.5pt}\end{center}

\subsection{4. Geometric Origin of Half-Integer
Spin}\label{geometric-origin-of-half-integer-spin}

\subsubsection{4.1 Edge Midpoints and Half-Integer
Coordinates}\label{edge-midpoints-and-half-integer-coordinates}

The vertices of the 5-dimensional orthoplex lie on the integer lattice
(combinations of \(\pm 1\) and \(0\) along each axis).

The midpoint of the edge connecting two adjacent vertices, for example
\(e_1 = (1,0,0,0,0)\) and \(e_2 = (0,1,0,0,0)\), is

\[m_{12} = \left(\frac{1}{2}, \frac{1}{2}, 0, 0, 0\right),\]

which has \textbf{half-integer coordinates}.

\textbf{Proposition 4.1.} In the 5-dimensional orthoplex defined on the
integer lattice, the edge midpoints lie on a half-integer lattice. These
half-integer lattice points constitute the geometric locus of
half-integer spin (fermion) states.

\subsubsection{4.2 Geometric Meaning of 720°
Rotation}\label{geometric-meaning-of-720-rotation}

A vertex state (boson) returns to the same vertex after one full
rotation (360°) on the lattice.

An edge-midpoint state (fermion) moves to an adjacent midpoint after one
rotation and returns to the original midpoint only after two full
rotations (720°).

This is the geometric expression of the property that a
spin-\(\frac{1}{2}\) particle acquires a phase factor of \((-1)\) under
360° rotation and returns to its original state only after 720°
rotation.

\subsubsection{4.3 Geometric Derivation of the Pauli Exclusion
Principle}\label{geometric-derivation-of-the-pauli-exclusion-principle}

\textbf{Proposition 4.2.} For half-integer spin states (edge midpoints),
the value \(s_z = 0\) does not exist. Therefore, the sign of the volume
is always determined to be either \(+\) or \(-\).

\emph{Proof.} The \(z\)-component of spin \(s\) takes values
\(s_z = -s, -s+1, \ldots, s-1, s\). When \(s\) is a half-integer,
\(s_z = 0\) does not appear in this sequence. \(\square\)

\textbf{Corollary 4.3.} In states where the volume sign is always
determined, two states with the same sign and the same quantum numbers
cannot coexist at the same spatial point. This corresponds to the Pauli
exclusion principle.

In contrast, for integer spin (bosons), the state \(s_z = 0\) exists and
admits both \(\pm\) volume signs, permitting multiple occupation of the
same state.

\begin{center}\rule{0.5\linewidth}{0.5pt}\end{center}

\subsection{5. Volume Sign and Force
Directionality}\label{volume-sign-and-force-directionality}

\subsubsection{5.1 Volume Signs of All Spin
States}\label{volume-signs-of-all-spin-states}

We classify each spin state \((s, s_z)\) by the sign of its volume.

\textbf{Bosonic states (integer spin):}

{\def\LTcaptype{none} % do not increment counter
\begin{longtable}[]{@{}ccc@{}}
\toprule\noalign{}
\(s\) & \(s_z\) & Volume sign \\
\midrule\noalign{}
\endhead
\bottomrule\noalign{}
\endlastfoot
0 & 0 & \(+\) \\
1 & \(+1\) & \(+\) \\
1 & \(0\) & \(\pm\) \\
1 & \(-1\) & \(-\) \\
2 & \(+2\) & \(+\) \\
2 & \(+1\) & \(+\) \\
2 & \(0\) & \(\pm\) \\
2 & \(-1\) & \(-\) \\
2 & \(-2\) & \(-\) \\
\end{longtable}
}

\textbf{Fermionic states (half-integer spin):}

{\def\LTcaptype{none} % do not increment counter
\begin{longtable}[]{@{}ccc@{}}
\toprule\noalign{}
\(s\) & \(s_z\) & Volume sign \\
\midrule\noalign{}
\endhead
\bottomrule\noalign{}
\endlastfoot
1/2 & \(+1/2\) & \(+\) \\
1/2 & \(-1/2\) & \(-\) \\
3/2 & \(+3/2\) & \(+\) \\
3/2 & \(+1/2\) & \(+\) \\
3/2 & \(-1/2\) & \(-\) \\
3/2 & \(-3/2\) & \(-\) \\
5/2 & \(+5/2\) & \(+\) \\
5/2 & \(+3/2\) & \(+\) \\
5/2 & \(+1/2\) & \(+\) \\
5/2 & \(-1/2\) & \(-\) \\
5/2 & \(-3/2\) & \(-\) \\
5/2 & \(-5/2\) & \(-\) \\
\end{longtable}
}

\subsubsection{5.2 Parity of Sign Reversals and Force
Directionality}\label{parity-of-sign-reversals-and-force-directionality}

\textbf{Proposition 5.1.} For bosons with an even number of sign
reversals (\(s = 0, 2\)), the product of volume signs is
\((-1)^n = +1\), and the force (if mediated) is unidirectional
(attractive only). For bosons with an odd number of sign reversals
(\(s = 1\)), the product is \((-1)^n = -1\), and the force is
bidirectional (both attractive and repulsive).

{\def\LTcaptype{none} % do not increment counter
\begin{longtable}[]{@{}cccc@{}}
\toprule\noalign{}
Sign reversals & Sign product & Bosonic spin & Force directionality \\
\midrule\noalign{}
\endhead
\bottomrule\noalign{}
\endlastfoot
0 (even) & \((+)(+) = +\) & 0 & None (scalar field) \\
1 (odd) & \((+)(-) = -\) & 1 & \(\pm\) (attractive and repulsive) \\
2 (even) & \((-)(-) = +\) & 2 & \(+\) only (attractive only) \\
\end{longtable}
}

\textbf{Remark.} The fact that the electromagnetic force (mediated by
spin-1 particles) exhibits both attraction and repulsion, while gravity
(mediated by spin-2 particles) is attractive only, is consistent with
experiment.

\subsubsection{5.3 Qualitative Difference Between the Signs of Spin 0
and Spin
2}\label{qualitative-difference-between-the-signs-of-spin-0-and-spin-2}

Spin 0 (zero sign reversals) and spin 2 (two sign reversals) both yield
a positive volume product, but their origins are fundamentally
different.

{\def\LTcaptype{none} % do not increment counter
\begin{longtable}[]{@{}
  >{\centering\arraybackslash}p{(\linewidth - 4\tabcolsep) * \real{0.1579}}
  >{\centering\arraybackslash}p{(\linewidth - 4\tabcolsep) * \real{0.4211}}
  >{\centering\arraybackslash}p{(\linewidth - 4\tabcolsep) * \real{0.4211}}@{}}
\toprule\noalign{}
\begin{minipage}[b]{\linewidth}\centering
\end{minipage} & \begin{minipage}[b]{\linewidth}\centering
Spin 0
\end{minipage} & \begin{minipage}[b]{\linewidth}\centering
Spin 2
\end{minipage} \\
\midrule\noalign{}
\endhead
\bottomrule\noalign{}
\endlastfoot
Volume product & \(+\) & \(+\) \\
Origin & \((+)(+)\): no reversal & \((-)(-)\): double reversal \\
Geometric property & Identity transformation & Self-inverse of
reversal \\
\end{longtable}
}

\textbf{Proposition 5.2.} The \((+)(+)\) type and the \((-)(-)\) type
both produce a positive volume product, but their algebraic structures
differ. The \((+)(+)\) type is the identity transformation, while the
\((-)(-)\) type is the self-inverse of sign reversal.

\textbf{Remark.} In particle physics, the Higgs particle (spin 0) is the
source of mass and is itself massive, whereas the graviton (spin 2) is
massless. This distinction corresponds to the algebraic difference
between the \((+)(+)\) type and the \((-)(-)\) type.

\begin{center}\rule{0.5\linewidth}{0.5pt}\end{center}

\subsection{6. Asymmetry of Volume
Signs}\label{asymmetry-of-volume-signs}

\subsubsection{6.1 Counts of Positive and Negative
States}\label{counts-of-positive-and-negative-states}

We classify all 21 spin states by their volume signs.

{\def\LTcaptype{none} % do not increment counter
\begin{longtable}[]{@{}cc@{}}
\toprule\noalign{}
Classification & Count \\
\midrule\noalign{}
\endhead
\bottomrule\noalign{}
\endlastfoot
Always \(+\) & 10 \\
Always \(-\) & 9 \\
Both \(\pm\) & 2 \\
\textbf{Total} & \textbf{21} \\
\end{longtable}
}

\textbf{Proposition 6.1.} In the orientation structure of the
5-dimensional orthoplex, the number of states with always-positive
volume (10) exceeds the number of states with always-negative volume (9)
by exactly one. This difference arises because the spin-0 state admits
only the \(+\) sign.

\subsubsection{6.2 The Number 10 and the Vertex
Count}\label{the-number-10-and-the-vertex-count}

The number of always-positive volume states, \textbf{10}, coincides with
the number of vertices of the 5-dimensional orthoplex.

\textbf{Remark.} The observational fact that the universe contains
slightly more matter than antimatter (baryon asymmetry) is one of the
outstanding problems in cosmology. The result of this paper suggests
that the origin of this asymmetry may lie in the sign structure of
spin-0 (scalar field) states.

\begin{center}\rule{0.5\linewidth}{0.5pt}\end{center}

\subsection{7. Classification of Edge Types and Particle
Hierarchy}\label{classification-of-edge-types-and-particle-hierarchy}

\subsubsection{\texorpdfstring{7.1 Edge Types and \(S_3\)
Symmetry}{7.1 Edge Types and S\_3 Symmetry}}\label{edge-types-and-s_3-symmetry}

Each edge of the 5-dimensional orthoplex connects two vertices on
different axes. The number of ways to choose 2 axes from 5 is
\(\binom{5}{2} = 10\).

Under the permutation symmetry \(S_3\) of the three spatial axes, the 10
edge types are classified into four groups.

{\def\LTcaptype{none} % do not increment counter
\begin{longtable}[]{@{}cccc@{}}
\toprule\noalign{}
Edge type & Axis pair & Count & \(S_3\) orbit \\
\midrule\noalign{}
\endhead
\bottomrule\noalign{}
\endlastfoot
Spatial--spatial & \((x_1, x_2), (x_2, x_3), (x_1, x_3)\) & 3 & 1
orbit \\
Spatial--fourth & \((x_1, x_4), (x_2, x_4), (x_3, x_4)\) & 3 & 1
orbit \\
Spatial--fifth & \((x_1, x_5), (x_2, x_5), (x_3, x_5)\) & 3 & 1 orbit \\
Fourth--fifth & \((x_4, x_5)\) & 1 & 1 orbit \\
\end{longtable}
}

\[10 = 3 + 3 + 3 + 1\]

\subsubsection{7.2 Meaning of the Threefold
Repetition}\label{meaning-of-the-threefold-repetition}

The \(S_3\) symmetry is the invariance under exchange of the three
spatial axes. Within each group, the three edge types differ only in the
spatial axis index \(i = 1, 2, 3\) and are geometrically equivalent.

\textbf{Proposition 7.1.} Under the \(S_3\) symmetry of the spatial
axes, the fermionic edge types form three groups of three.

\textbf{Remark.} The three-generation structure of particle physics
(e.g., \(e, \mu, \tau\)) is one of the outstanding unsolved problems.
The result of this paper suggests that the three generations may have a
geometric origin in the permutation symmetry of the three spatial axes.

\subsubsection{7.3 Correspondence with the Total Edge
Count}\label{correspondence-with-the-total-edge-count}

The total number of edges of the 5-dimensional orthoplex is 40. This can
be decomposed as follows.

\[40 = 10 \times 2 \times 2\]

{\def\LTcaptype{none} % do not increment counter
\begin{longtable}[]{@{}cc@{}}
\toprule\noalign{}
Factor & Geometric meaning \\
\midrule\noalign{}
\endhead
\bottomrule\noalign{}
\endlastfoot
10 & Edge types (\(\binom{5}{2}\)) \\
\(\times 2\) & Edge direction (\(\pm\)) \\
\(\times 2\) & Hemisphere selection (north/south) \\
\end{longtable}
}

\begin{center}\rule{0.5\linewidth}{0.5pt}\end{center}

\subsection{8. Relation to Prior Work}\label{relation-to-prior-work}

\subsubsection{8.1 Direct Precedents}\label{direct-precedents}

To the author's knowledge, no prior work derives spin values from the
orientation structure of the 5-dimensional orthoplex.

\subsubsection{8.2 Related Research}\label{related-research}

The following works are related to individual aspects of this paper.

\textbf{Topological derivation of spin-statistics:} - Finkelstein, D. \&
Rubinstein, J. (1968). ``Connection between Spin, Statistics, and
Kinks.'' \emph{J. Math. Phys.} 9, 1762.

\textbf{Geometric models of matter:} - Atiyah, M.F., Manton, N.S. \&
Schroers, B.J. (2012). ``Geometric Models of Matter.'' \emph{Proc. R.
Soc. A} 468, 1252. arXiv: 1108.5151.

\textbf{Clifford algebras and the Standard Model:} - Furey, C. (2016).
``Standard Model Physics from an Algebra?'' arXiv: 1611.09182.

\textbf{5-dimensional unification:} - Kaluza, Th. (1921). ``Zum
Unitätsproblem der Physik.'' \emph{Sitzungsber. Preuss. Akad. Wiss.
Berlin} 966--972.

\subsubsection{8.3 Numerical
Correspondences}\label{numerical-correspondences}

We record the correspondences between the numbers arising in this paper
and known mathematical structures.

{\def\LTcaptype{none} % do not increment counter
\begin{longtable}[]{@{}cc@{}}
\toprule\noalign{}
Number in this paper & Known structure \\
\midrule\noalign{}
\endhead
\bottomrule\noalign{}
\endlastfoot
21 (total spin states) & \(\dim \mathrm{SO}(7) = 21\) \\
32 (4-dimensional facets) & \(\dim \mathrm{Cl}(5) = 2^5 = 32\) \\
10 (vertices) & \(\dim \mathrm{SO}(5) = 10\) \\
\end{longtable}
}

Whether these coincidences are accidental or fundamental is left as a
question for future work.

\begin{center}\rule{0.5\linewidth}{0.5pt}\end{center}

\subsection{9. Open Problems and Future
Directions}\label{open-problems-and-future-directions}

The results of this paper are self-contained as propositions of discrete
geometry, but the following problems remain open.

\subsubsection{9.1 Quantitative Derivation of Mass
Ratios}\label{quantitative-derivation-of-mass-ratios}

We suggested that the mass hierarchy of three generations may arise from
the breaking of the \(S_3\) symmetry of spatial axes by the choice of
projection axis (Section 7). However, the specific values of mass ratios
have not been derived.

\textbf{Problem:} Can the mass ratios \(m_e : m_\mu : m_\tau\), etc., be
derived quantitatively from the angles between the projection axis and
each spatial axis?

\subsubsection{\texorpdfstring{9.2 Derivation of the Gauge Group
\(\mathrm{SU}(3) \times \mathrm{SU}(2) \times \mathrm{U}(1)\)}{9.2 Derivation of the Gauge Group \textbackslash mathrm\{SU\}(3) \textbackslash times \textbackslash mathrm\{SU\}(2) \textbackslash times \textbackslash mathrm\{U\}(1)}}\label{derivation-of-the-gauge-group-mathrmsu3-times-mathrmsu2-times-mathrmu1}

The permutation group \(S_3\) of the spatial axes is related to
\(\mathrm{SU}(3)\) as a Weyl group, but whether the full Standard Model
gauge group can be derived from the symmetry of the 5-dimensional
orthoplex remains unproven.

\textbf{Problem:} Can the correspondences \(S_3 \to \mathrm{SU}(3)\) and
edge-type structure \(\to \mathrm{SU}(2) \times \mathrm{U}(1)\) be
rigorously constructed?

\subsubsection{9.3 Discrepancy Between 40 Edges and 30 Standard Model
States}\label{discrepancy-between-40-edges-and-30-standard-model-states}

The number of fundamental fermionic states per generation in the
Standard Model is 15 (or 30 including antiparticles), whereas the
decomposition of the 5-dimensional orthoplex edges as type \(\times\)
direction \(\times\) hemisphere yields 40. The physical meaning of the
difference of 10 is unclear.

\textbf{Problem:} Does the discrepancy between 40 and 30 predict the
existence of undiscovered states such as right-handed neutrinos, or does
the counting correspondence require modification?

\subsubsection{\texorpdfstring{9.4 Meaning of
\(21 = \dim \mathrm{SO}(7)\)}{9.4 Meaning of 21 = \textbackslash dim \textbackslash mathrm\{SO\}(7)}}\label{meaning-of-21-dim-mathrmso7}

The coincidence between the total number of spin states (21) and the
dimension of \(\mathrm{SO}(7)\) is unexplained.

\textbf{Problem:} Is the spin state space of the 5-dimensional orthoplex
isomorphic to the adjoint representation of \(\mathrm{SO}(7)\) or a
subspace thereof?

\subsubsection{9.5 Geometric Origin of Coupling
Constants}\label{geometric-origin-of-coupling-constants}

This paper derived spin values and force directionality but did not
address the magnitudes of couplings (coupling constants).

\textbf{Problem:} Do values such as the electromagnetic coupling
constant \(\alpha \approx 1/137\) correspond to geometric quantities
(angles, area ratios, etc.) of the 5-dimensional orthoplex?

\begin{center}\rule{0.5\linewidth}{0.5pt}\end{center}

\subsection{10. Conclusion}\label{conclusion}

The following results have been derived purely geometrically from the
orientation structure of the 5-dimensional orthoplex.

\begin{enumerate}
\def\labelenumi{\arabic{enumi}.}
\tightlist
\item
  \textbf{Complete classification of spin values}: The six values
  \(s = 0, \frac{1}{2}, 1, \frac{3}{2}, 2, \frac{5}{2}\) are necessarily
  determined by the sign-reversal structure of two non-spatial axes and
  the half-integer structure of edge midpoints in 5 dimensions.
\item
  \textbf{Derivation of the fermion--boson distinction}: The existence
  or non-existence of the \(s_z = 0\) state provides a geometric account
  of the differing statistical properties of integer and half-integer
  spins.
\item
  \textbf{Force directionality}: The parity of the number of sign
  reversals determines the directionality of the mediated force
  (unidirectional versus bidirectional).
\item
  \textbf{Origin of particle hierarchy}: The \(3 + 3 + 3 + 1\)
  decomposition of edge types under \(S_3\) symmetry provides a
  geometric candidate for the generation structure.
\item
  \textbf{Matter--antimatter asymmetry}: The difference of 1 between 10
  positive-volume states and 9 negative-volume states is attributable to
  the sign structure of spin 0.
\end{enumerate}

These results indicate the possibility that the axioms of quantum
mechanics concerning spin (the spin-statistics theorem, the Pauli
exclusion principle) may be re-derived as theorems of 5-dimensional
integer geometry.

\begin{center}\rule{0.5\linewidth}{0.5pt}\end{center}

\subsection{References}\label{references}

{[}1{]} N. Kihara, ``Geometric Formulation of 4-Dimensional Space via
Central Projection,'' 2026. DOI: 10.5281/zenodo.19427780

{[}2{]} N. Kihara, ``Geometric Symmetry of Central Projection:
Mathematical Foundations of the Multi-Axis Model,'' 2026. DOI:
10.5281/zenodo.19434932

{[}3{]} N. Kihara, ``Formulation of Structural Requirements for a Theory
of Everything,'' 2026. DOI: 10.5281/zenodo.19601592

{[}4{]} N. Kihara, ``Classification and Structural Analysis of the 19
Arbitrary Parameters of the Standard Model,'' 2026. DOI:
10.5281/zenodo.19604965

{[}5{]} N. Kihara, ``Full Spherical Coverage of Central Projection in
Discrete Space and the Number-Theoretic Necessity of 5-Dimensional
Background Space,'' 2026. DOI: 10.5281/zenodo.19624957

{[}6{]} D. Finkelstein \& J. Rubinstein, ``Connection between Spin,
Statistics, and Kinks,'' \emph{J. Math. Phys.} 9, 1762 (1968).

{[}7{]} M.F. Atiyah, N.S. Manton \& B.J. Schroers, ``Geometric Models of
Matter,'' \emph{Proc. R. Soc. A} 468, 1252 (2012). arXiv: 1108.5151.

{[}8{]} C. Furey, ``Standard Model Physics from an Algebra?'' arXiv:
1611.09182 (2016).

{[}9{]} Th. Kaluza, ``Zum Unitätsproblem der Physik,''
\emph{Sitzungsber. Preuss. Akad. Wiss. Berlin} 966--972 (1921).

{[}10{]} H.S.M. Coxeter, \emph{Regular Polytopes}, Dover (1973).

\end{document}
